\subsection{Report-driven evidence retrieval}
\label{sec:rag-diag}

The condition report serves two audiences. Its chronological narrative and optional JSON snapshots support human inspection and audit, whereas its compact alert summary is the machine-facing interface to retrieval. The summary contains variable-wise alerts, spike and zero-event counts, morphology, collinearity, lag, and selected spatial-snapshot cues. Queries are constructed from this structured representation rather than from a short operator phrase or an unreviewed free-text diagnosis hypothesis.

Source PDFs are page-cropped when required, converted to Markdown, chunked, and embedded in a vector store~\citep{liu2023llamaindex}. At retrieval time, the system first searches for passages relevant to the observed phenomenon or mechanism, then may supplement them with passages on redundancy, topology, monitoring practice, or recommended checks. A memo composer links candidate interpretations to the retrieved passages with numbered citations. It must not convert this support into an automatically confirmed root cause.

This design makes the report the semantic bridge between the measurement and evidence resources introduced in \cref{sec:data}. It also explains why the retrieval experiments in \cref{sec:eval} are part of the system contribution: they stress-test whether this evidence layer preserves the intended source policy under controlled corpus expansion, query variation, and chunking.
